\section{总结与延伸}

本讲把儿童在线安全从单一内容分类问题展开为完整 frontier system。最重要的不是某个产品名，而是以下可迁移原则：

\begin{enumerate}
  \item \textbf{端到端结果优先}：signal、platform decision、lawful report 与 victim protection 是不同阶段，必须有清晰 handoff；
  \item \textbf{Known 与 unknown 分层}：hash matching 高速处理已知对象，predictive detection 扩大未知候选覆盖，两者都需要 policy 和审核；
  \item \textbf{模型与队列共同设计}：precision、recall、calibration、threshold、review capacity 与 moderator wellbeing 不能拆开验收；
  \item \textbf{数据治理是核心能力}：trusted source、provenance、segregation、retention、label protocol 和 audit 决定模型是否可部署；
  \item \textbf{Safety by Design 贯穿生命周期}：development、deployment、maintenance、red teaming 与 transparency 都有明确 owner；
  \item \textbf{隐私是架构目标}：E2EE、metadata、endpoint 和 network signal 各有权利与安全边界，必须最小化并可申诉；
  \item \textbf{更早信号需要更克制动作}：conversation 与 graph analysis 提供早期预警，但动作应分层、可逆并保留 human review；
  \item \textbf{生态角色不可互换}：平台、技术提供者、报告机构、调查人员与支持服务通过 contract 和 audit 形成保护闭环。
\end{enumerate}

\begin{importantbox}{平台安全作业：设计一个可验收闭环}
选择一个允许用户上传内容或私信的产品，提交一页设计：画出检测、复核、处置与报告状态机；定义 known-match 与 predictive signal 的字段；给出 precision/recall 与 reviewer-capacity 约束；写出 privacy、retention、appeal、incident 和 provider-outage 策略。禁止只画“AI moderation”方框，必须标明每个 handoff 的 owner 与证据。
\end{importantbox}

\subsection{拓展阅读}

{\small
\begin{itemize}[nosep,leftmargin=*]
  \item \href{https://www.thorn.org/solutions/for-platforms/}{\textbf{Thorn platform solutions}}：Safer detection、Match 与 Predict 的官方入口。
  \item \href{https://info.thorn.org/hubfs/thorn-safety-by-design-for-generative-AI.pdf}{\textbf{Thorn Safety by Design for Generative AI}}：生成式 AI 生命周期原则白皮书。
  \item \href{https://www.missingkids.org/gethelpnow/cybertipline/cybertiplinedata}{\textbf{NCMEC CyberTipline data}}：年度报告口径与趋势说明。
  \item \href{https://www.missingkids.org/netsmartz/home}{\textbf{NCMEC NetSmartz}}：面向家庭、学校和儿童的预防教育资源。
  \item \href{https://nvlpubs.nist.gov/nistpubs/ai/NIST.AI.600-1.pdf}{\textbf{NIST AI 600-1}}：Generative AI Profile。
  \item \href{https://nvlpubs.nist.gov/nistpubs/ai/NIST.AI.100-4.pdf}{\textbf{NIST AI 100-4}}：synthetic-content risk technical approaches。
\end{itemize}
}

\end{document}
